\documentclass[12pt,a4paper]{article}

% ------------------------------------------------
% PACOTES
% ------------------------------------------------
\usepackage[utf8]{inputenc}
\usepackage[T1]{fontenc}
\usepackage[brazil]{babel}
\usepackage{geometry}
\usepackage{setspace}
\usepackage{hyperref}
\usepackage{graphicx}
\usepackage{booktabs}

% ------------------------------------------------
% LAYOUT
% ------------------------------------------------
\geometry{
  left=3cm,
  right=2cm,
  top=3cm,
  bottom=2.5cm
}

\setstretch{1.5}

% ------------------------------------------------
\begin{document}

\begin{center}
    \Large\textbf{RELATÓRIO TÉCNICO}\\[0.3cm]
    \large\textbf{Projeto-Piloto do Sistema de Presença com QR Code Dinâmico}\\[0.3cm]
    \normalsize Faculdade de Engenharia Civil -- UNICAMP\\
    Disciplina: \textbf{Fundações}
\end{center}

\vspace{1cm}

% ------------------------------------------------

\section*{1. Introdução}

Este relatório apresenta os resultados do projeto-piloto de implementação de um sistema automatizado de controle de presença em sala de aula com utilização de QR Codes dinâmicos, aplicado à disciplina \textbf{Fundações}, envolvendo aproximadamente 100 estudantes.

% ------------------------------------------------

\section*{2. Objetivos}

\begin{itemize}
    \item Avaliar a viabilidade operacional do sistema;
    \item Medir o impacto no tempo de chamada;
    \item Verificar confiabilidade e segurança dos registros;
    \item Analisar o grau de aceitação discente e docente.
\end{itemize}

% ------------------------------------------------

\section*{3. Metodologia}

\subsection*{3.1 Ambiente}

O sistema foi hospedado em plataforma em nuvem (Render/AWS), com backend desenvolvido em Python (FastAPI) e banco de dados em planilha Google Sheets integrada via API.

\subsection*{3.2 Procedimento}

Em cada aula foi projetado um QR Code temporário válido por até 3 minutos, permitindo que os alunos realizassem o registro de presença via smartphone mediante inserção de RA e nome.

\subsection*{3.3 Segurança}

Foram aplicados:

\begin{itemize}
    \item Token exclusivo por aula;
    \item Expiração automática;
    \item Bloqueio de duplicidades por RA;
    \item Registro de horário do check-in.
\end{itemize}

% ------------------------------------------------

\section*{4. Resultados}

\subsection*{4.1 Indicadores Quantitativos}

Inserir tabela resumo:

\begin{center}
\begin{tabular}{lcc}
\toprule
\textbf{Indicador} & \textbf{Chamada Manual} & \textbf{QR Code} \\
\midrule
Tempo médio (min) & --- & --- \\
Taxa de registro (\%) & --- & --- \\
Falhas reportadas (\%) & --- & --- \\
\bottomrule
\end{tabular}
\end{center}

\subsection*{4.2 Indicadores Qualitativos}

Relatar:

\begin{itemize}
    \item Percepção dos alunos;
    \item Feedback dos docentes;
    \item Facilidade operacional.
\end{itemize}

% ------------------------------------------------

\section*{5. Discussão}

Comparar o desempenho do método tradicional com o sistema digital, destacando ganhos operacionais, limitações observadas e eventuais ajustes necessários.

% ------------------------------------------------

\section*{6. Conformidade com LGPD}

O sistema coletou apenas dados mínimos para fins de controle acadêmico, em conformidade com a legislação vigente, sem coleta de dados sensíveis.

% ------------------------------------------------

\section*{7. Conclusão}

Avaliar a viabilidade de expansão do sistema para outras disciplinas da FEC ou integração futura aos sistemas institucionais.

% ------------------------------------------------

\section*{8. Encaminhamentos}

Recomendações:

\begin{itemize}
    \item Ampliação do piloto;
    \item Implementação de autenticação institucional;
    \item Criação de painéis de acompanhamento.
\end{itemize}

% ------------------------------------------------

\section*{9. Coordenação do Projeto}

\begin{center}
\textbf{Heloi Moacyr Tanoue}
\end{center}

% ------------------------------------------------

\vfill
\begin{center}
    Novembro de 2025
\end{center}

\end{document}

